\chapter{Zupełność jednostajna i uzupełnienia}
\label{ch:topog-uzupelnienia}

\noindent\textbf{Cel rozdziału.}
W rozdziale~12 wprowadziliśmy sieci i filtry Cauchy'ego,
a w rozdziale~11 zupełność przestrzeni metrycznych.
Wyjaśnimy ich związek, skonstruujemy uzupełnienie
przestrzeni jednostajnej rozdzielającej oraz
udowodnimy jednoznaczność i własność uniwersalną tej konstrukcji.
Oddzielimy zupełność ustalonej struktury jednostajnej
od zupełnej metryzowalności przestrzeni topologicznej.

\section{Zupełność w języku filtrów i sieci}
\label{sec:topog-uzupelnienia-definicje}

\begin{definicja}[Zbieżność filtra i zupełność jednostajna]
\label{def:topog-uzupelnienia-zupelnosc}
Niech $(X,\mathcal U)$ będzie przestrzenią jednostajną.
Filtr $\mathcal F$ na $X$ jest \emph{zbieżny do $x\in X$},
gdy każde otoczenie $x$ w topologii $\tau_{\mathcal U}$
należy do $\mathcal F$.
Przestrzeń $(X,\mathcal U)$ jest \emph{zupełna},
jeżeli każdy filtr Cauchy'ego na $X$ jest zbieżny
do co najmniej jednego punktu $X$.
\end{definicja}

\begin{uwaga}[Rozdzielanie i przestrzeń pusta]
Bez założenia Hausdorffa granica filtra nie musi być
jednoznaczna. Na przestrzeni pustej nie ma filtrów
właściwych, zatem jest ona zupełna.
\end{uwaga}

\begin{twierdzenie}[Równoważność filtrów i sieci Cauchy'ego]
\label{thm:topog-uzupelnienia-filtry-sieci}
Przestrzeń jednostajna $(X,\mathcal U)$ jest zupełna
wtedy i tylko wtedy, gdy każda jej sieć Cauchy'ego
jest zbieżna w topologii $\tau_{\mathcal U}$.
\end{twierdzenie}
\begin{proof}
Załóżmy, że wszystkie filtry Cauchy'ego są zbieżne.
Jeśli $(x_\alpha)_{\alpha\in D}$ jest siecią Cauchy'ego,
to jej filtr ogonów jest Cauchy'ego na mocy
twierdzenia o filtrze ogonów z rozdziału~12.
Niech zbiega do $x$. Dla każdego otoczenia $O$
punktu $x$ istnieje ogon zawarty w $O$,
więc $x_\alpha\to x$.

Załóżmy teraz, że każda sieć Cauchy'ego jest zbieżna.
Niech $\mathcal F$ będzie filtrem Cauchy'ego.
Rozważmy zbiór skierowany
$D=\{(A,a):A\in\mathcal F,\ a\in A\}$,
z porządkiem $(A,a)\preceq(B,b)$ wtedy i tylko wtedy,
gdy $B\subseteq A$. Jest on niepusty, ponieważ
filtry są właściwe, i skierowany: dla
$(A,a),(B,b)$ wybieramy $c\in A\cap B$
oraz $(A\cap B,c)$. Definiujemy sieć
$x_{(A,a)}=a$.
Dla dowolnego $E\in\mathcal U$ istnieje
$A_0\in\mathcal F$ z $A_0\times A_0\subseteq E$.
Od indeksu $(A_0,a_0)$ wszystkie wyrazy
sieci leżą w $A_0$, więc sieć jest Cauchy'ego.
Niech zbiega do $x$.
Jeżeli $O$ jest otoczeniem $x$, to istnieje
$(A,a)\in D$ takie, że wszystkie późniejsze
wyrazy leżą w $O$. Dla każdego $b\in A$
indeks $(A,b)$ jest późniejszy od $(A,a)$,
a zatem $b\in O$. Stąd $A\subseteq O$
i $O\in\mathcal F$. Filtr $\mathcal F$ zbiega do $x$.
\end{proof}

\begin{stwierdzenie}[Zupełność metryczna i jednostajna]
\label{prop:topog-uzupelnienia-metryczna}
Niech $(X,d)$ będzie przestrzenią metryczną,
a $\mathcal U_d$ strukturą jednostajną wyznaczoną
przez $d$. Wtedy $(X,d)$ jest zupełna wtedy
i tylko wtedy, gdy $(X,\mathcal U_d)$ jest zupełna.
\end{stwierdzenie}
\begin{proof}
Jeżeli $(X,\mathcal U_d)$ jest zupełna, to
każdy ciąg Cauchy'ego jest siecią Cauchy'ego
w $\mathcal U_d$, zatem zbiega w topologii metrycznej.

Odwrotnie, niech $(X,d)$ będzie zupełna, a
$\mathcal F$ filtrem Cauchy'ego. Dla każdego
$n\ge1$ wybieramy $A_n\in\mathcal F$ takie, że
$d(a,b)<2^{-n}$ dla wszystkich $a,b\in A_n$.
Zastępujemy $A_n$ przez $\bigcap_{k=1}^n A_k$;
wówczas $A_{n+1}\subseteq A_n$ i $A_n\ne\varnothing$.
Wybierzmy $x_n\in A_n$. Gdy $m\ge n$,
punkty $x_m,x_n$ należą do $A_n$,
więc $d(x_m,x_n)<2^{-n}$.
Z zupełności $x_n\to x\in X$.
Dla $\varepsilon>0$ wybieramy $n$ tak,
że $2^{-n}<\varepsilon/2$ i
$d(x_n,x)<\varepsilon/2$.
Wtedy $A_n\subseteq B_d(x,\varepsilon)$,
więc kula ta należy do $\mathcal F$.
To dowodzi zbieżności filtra do $x$.
\end{proof}

\begin{przyklad}[Jednakowa topologia, inna zupełność]
\label{ex:topog-uzupelnienia-arctan}
Na $\R$ metryki $d(x,y)=|x-y|$ oraz
$e(x,y)=|\arctan x-\arctan y|$ wyznaczają tę
samą topologię (rozdział~12).
Przestrzeń $(\R,d)$ jest zupełna, lecz
$(\R,e)$ nie jest: ciąg $x_n=n$ jest
Cauchy'ego dla $e$, natomiast nie ma granicy
w zwykłej topologii $\R$.
Zupełność dotyczy struktury jednostajnej,
a nie samej topologii.
\end{przyklad}

\section{Uzupełnienie przestrzeni metrycznej}
\label{sec:topog-uzupelnienia-metryczne}

\begin{definicja}[Uzupełnienie metryczne]
\label{def:topog-uzupelnienia-metryczne}
\emph{Uzupełnieniem metrycznym} przestrzeni
metrycznej $(X,d)$ nazywamy parę
$(\widehat X,\iota)$, gdzie
$(\widehat X,\widehat d)$ jest zupełną przestrzenią
metryczną, a $\iota:(X,d)\to(\widehat X,\widehat d)$
jest izometrią o gęstym obrazie.
\end{definicja}

\begin{twierdzenie}[Konstrukcja uzupełnienia metrycznego]
\label{thm:topog-uzupelnienia-metryczne}
Każda przestrzeń metryczna $(X,d)$ ma uzupełnienie
metryczne. Jeżeli $X\ne\varnothing$,
można je zbudować z klas ciągów Cauchy'ego w $X$.
\end{twierdzenie}
\begin{proof}
Dla $X=\varnothing$ przyjmujemy $\widehat X=\varnothing$
i jedyne możliwe odwzorowanie $\iota$.
Niech dalej $X\ne\varnothing$. Oznaczmy przez
$\mathcal C(X)$ zbiór ciągów Cauchy'ego w $X$
i połóżmy $(x_n)\sim(y_n)$, gdy
$d(x_n,y_n)\to0$. Zwrotność i symetria są
bezpośrednie, a przechodniość wynika z
$d(x_n,z_n)\le d(x_n,y_n)+d(y_n,z_n)$.
Definiujemy $\widehat X=\mathcal C(X)/{\sim}$.

Dla ciągów Cauchy'ego $(x_n),(y_n)$
ciąg liczb $d(x_n,y_n)$ jest Cauchy'ego, gdyż
\[
 |d(x_n,y_n)-d(x_m,y_m)|
 \le d(x_n,x_m)+d(y_n,y_m).
\]
Istnieje zatem granica; definiujemy
$\widehat d([(x_n)],[(y_n)])=\lim_nd(x_n,y_n)$.
Zmiana przedstawicieli nie zmienia granicy,
ponieważ różnica odległości jest nie większa
niż suma odległości między odpowiednimi
przedstawicielami. Symetria i nierówność
trójkąta przechodzą na granicę.
Jeżeli $\widehat d=0$, ciągi są równoważne,
więc $\widehat d$ jest metryką.

Odwzorowanie $\iota(x)=[(x,x,\ldots)]$
jest izometrią. Dla $z=[(x_n)]$ mamy
\[
 \widehat d(\iota(x_n),z)=\lim_{m\to\infty}d(x_n,x_m)
 \longrightarrow0.
\]
Ostatnie przejście wynika wprost z warunku
Cauchy'ego: dla dużego $n$ wszystkie
$d(x_n,x_m)$ z $m$ dostatecznie dużym
są mniejsze od zadanego $\varepsilon$.
Stąd $\iota(X)$ jest gęsty w $\widehat X$.

Niech $(z_k)$ będzie ciągiem Cauchy'ego
w $(\widehat X,\widehat d)$.
Z gęstości wybieramy $x_k\in X$ z
$\widehat d(z_k,\iota(x_k))<2^{-k}$.
Wtedy
\[
 d(x_k,x_l)\le
 2^{-k}+\widehat d(z_k,z_l)+2^{-l},
\]
więc $(x_k)$ jest ciągiem Cauchy'ego.
Niech $z=[(x_k)]$. Z udowodnionej gęstości
$\iota(x_k)\to z$; nierówność
$\widehat d(z_k,z)\le
\widehat d(z_k,\iota(x_k))+
\widehat d(\iota(x_k),z)$
daje $z_k\to z$. Uzupełnienie jest zupełne.
\end{proof}

\begin{przyklad}[Uzupełnienie liczb wymiernych]
\label{ex:topog-uzupelnienia-q}
Zbiór $\Q$ ze zwykłą metryką jest gęsty w
zupełnej przestrzeni $\R$. Dlatego
$(\R,|\cdot|)$ jest uzupełnieniem metrycznym
$(\Q,|\cdot|)$. Klasa ciągu wymiernych
przybliżeń liczby $\sqrt2$ odpowiada
punktowi $\sqrt2\in\R$.
\end{przyklad}

\section{Filtry Cauchy'ego i minimalność}
\label{sec:topog-uzupelnienia-filtry}

\begin{definicja}[Filtry bliskie]
\label{def:topog-uzupelnienia-bliskie}
Niech $(X,\mathcal U)$ będzie przestrzenią
jednostajną. Dwa filtry Cauchy'ego
$\mathcal F,\mathcal G$ na $X$ nazywamy
\emph{bliskimi}, i piszemy $\mathcal F\sim\mathcal G$,
gdy dla każdego $E\in\mathcal U$
istnieją $A\in\mathcal F$ i $B\in\mathcal G$
takie, że $A\times B\subseteq E$.
\end{definicja}

\begin{lemat}[Relacja bliskości jest równoważnością]
\label{lem:topog-uzupelnienia-rownowaznosc}
Relacja $\sim$ jest relacją równoważności
na zbiorze filtrów Cauchy'ego w $X$.
\end{lemat}
\begin{proof}
Zwrotność wynika z warunku Cauchy'ego:
dla każdego $E$ istnieje $A\in\mathcal F$
z $A\times A\subseteq E$.
Symetria: dla $E$ wybieramy $E\cap E^{-1}$
i zamieniamy role obu filtrów.

Dla przechodniości wybierzmy $E\in\mathcal U$
i symetryczne $D\in\mathcal U$ z
$D\circ D\subseteq E$. Jeżeli
$\mathcal F\sim\mathcal G$ oraz
$\mathcal G\sim\mathcal H$, to istnieją
$A\in\mathcal F$, $B,B'\in\mathcal G$,
$C\in\mathcal H$ z $A\times B\subseteq D$
i $B'\times C\subseteq D$.
Z właściwości filtra wybieramy
$b\in B\cap B'$. Dla wszystkich
$a\in A,c\in C$ mamy $(a,b),(b,c)\in D$,
więc $(a,c)\in D\circ D\subseteq E$.
Stąd $\mathcal F\sim\mathcal H$.
\end{proof}

\begin{definicja}[Filtr Cauchy'ego minimalny]
\label{def:topog-uzupelnienia-minimalny}
Filtr Cauchy'ego $\mathcal F$ jest
\emph{minimalny}, gdy nie istnieje
ściśle mniejszy w sensie inkluzji
filtr Cauchy'ego $\mathcal G\subsetneq\mathcal F$.
\end{definicja}

\begin{lemat}[Minimalny przedstawiciel]
\label{lem:topog-uzupelnienia-minimalizacja}
Każdy filtr Cauchy'ego $\mathcal F$ ma
dokładnie jeden minimalny filtr Cauchy'ego
bliski $\mathcal F$. Jest nim filtr
$\mathcal F^\circ$ generowany przez zbiory
$E[A]=\bigcup_{a\in A}E[a]$, gdzie
$E\in\mathcal U$ i $A\in\mathcal F$.
\end{lemat}
\begin{proof}
Zbiory $E[A]$ są niepuste. Ich rodzina
jest bazą filtra, ponieważ
$(E\cap D)[A\cap B]\subseteq E[A]\cap D[B]$.
Ponadto $A\subseteq E[A]$, zatem
$\mathcal F^\circ\subseteq\mathcal F$.

Ustalmy $V\in\mathcal U$. Wybierzmy
symetryczne $D\in\mathcal U$ z
$D\circ D\circ D\subseteq V$:
najpierw wybieramy otoczenie o kwadracie
zawartym w $V$, potem symetryczne otoczenie
o kwadracie zawartym w poprzednim.
Z warunku Cauchy'ego istnieje
$A\in\mathcal F$ z $A\times A\subseteq D$.
Jeżeli $u,v\in D[A]$, wybieramy
$a,b\in A$ z $(a,u),(b,v)\in D$.
Z symetrii $(u,a),(a,b),(b,v)\in D$,
stąd $(u,v)\in D^3\subseteq V$.
Zatem $\mathcal F^\circ$ jest Cauchy'ego.

Aby wykazać bliskość $\mathcal F^\circ$
i $\mathcal F$, dla $V$ wybieramy
symetryczne $D$ z $D^2\subseteq V$
oraz $A\in\mathcal F$ z $A\times A\subseteq D$.
Wtedy $D[A]\in\mathcal F^\circ$
i $D[A]\times A\subseteq D^2\subseteq V$.

Niech teraz $\mathcal G$ będzie dowolnym
filtrem Cauchy'ego bliskim $\mathcal F$.
Pokażemy $\mathcal F^\circ\subseteq\mathcal G$.
Dla bazowego $E[A]$, gdzie
$A\in\mathcal F$, wybieramy symetryczne
$D\subseteq E$. Bliskość zapewnia
$A'\in\mathcal F$ i $B\in\mathcal G$
z $A'\times B\subseteq D$.
Biorąc $a\in A\cap A'$ otrzymujemy
$B\subseteq D[a]\subseteq E[A]$,
więc $E[A]\in\mathcal G$.
Wobec tego $\mathcal F^\circ$
zawiera się w każdym bliskim filtrze Cauchy'ego.

W szczególności, jeśli $\mathcal H$
jest filtrem Cauchy'ego zawartym w
$\mathcal F^\circ$, to jest bliski
$\mathcal F^\circ$: dla $V$ weźmy
$A\in\mathcal H$ z $A\times A\subseteq V$;
ponieważ $\mathcal H\subseteq\mathcal F^\circ$,
mamy również $A\in\mathcal F^\circ$.
Z już wykazanej minimalności względem
filtrów bliskich wynika
$\mathcal F^\circ\subseteq\mathcal H$.
Stąd równość, czyli minimalność.
Jeśli dwa minimalne filtry są bliskie,
oba muszą się zawierać jeden w drugim,
więc są równe.
\end{proof}

\section{Konstrukcja uzupełnienia jednostajnego}
\label{sec:topog-uzupelnienia-konstrukcja}

\begin{twierdzenie}[Istnienie uzupełnienia jednostajnego]
\label{thm:topog-uzupelnienia-istnienie}
Niech $(X,\mathcal U)$ będzie rozdzielającą
przestrzenią jednostajną. Istnieje zupełna,
rozdzielająca przestrzeń jednostajna
$(\widehat X,\widehat{\mathcal U})$
oraz jednostajnie ciągłe odwzorowanie
$\iota:X\to\widehat X$ będące
izomorfizmem jednostajnym na swój gęsty obraz.
\end{twierdzenie}
\begin{proof}
Niech $\widehat X$ będzie zbiorem wszystkich
minimalnych filtrów Cauchy'ego na $X$.
Dla $E\in\mathcal U$ definiujemy relację
\[
 \widehat E=\{(\mathcal F,\mathcal G)\in
 \widehat X^2:\exists A\in\mathcal F,\ B\in\mathcal G,\
 A\times B\subseteq E\}.
\]
Relacje $\widehat E$ tworzą bazę struktury jednostajnej.
Istotnie, każda zawiera przekątną przez warunek
Cauchy'ego. Dla przecięć: wybierzmy symetryczne
$D$ z $D^3\subseteq E\cap F$.
Jeśli $(\mathcal A,\mathcal B)\in\widehat D$,
to należy także do $\widehat E\cap\widehat F$.
Odwrotność $\widehat E$ jest równa
$\widehat{E^{-1}}$. Gdy $D$ jest symetryczne
i $D^2\subseteq E$, mamy
$\widehat D\circ\widehat D\subseteq\widehat E$:
z $A\times B\subseteq D$ oraz
$B'\times C\subseteq D$ i
$b\in B\cap B'$ otrzymujemy
$A\times C\subseteq D^2\subseteq E$.
Domykanie relacji w górę daje
$\widehat{\mathcal U}$.

Jeśli $(\mathcal F,\mathcal G)$ należy do
wszystkich $\widehat E$, to
$\mathcal F\sim\mathcal G$, a więc
$\mathcal F=\mathcal G$ na mocy
lematu~\ref{lem:topog-uzupelnienia-minimalizacja}.
Struktura jest rozdzielająca.

Dla $x\in X$ filtr otoczeń $\mathcal N_x$
w $\tau_{\mathcal U}$ jest filtrem Cauchy'ego:
dla każdego $E$ wybieramy symetryczne $D$
z $D^2\subseteq E$, a
$D[x]\times D[x]\subseteq E$.
Ponadto jest minimalny: dla dowolnego
$A\in\mathcal N_x$ wybieramy $E$ z
$E[x]\subseteq A$. Ponieważ
$\{x\}\in\mathcal N_x$ nie musi zachodzić,
zamiast tego z twierdzenia o topologii
indukowanej bierzemy mniejsze $D[x]$
takie, że $D^2[x]\subseteq A$.
Wówczas $D[D[x]]\subseteq A$, więc
$A\in\mathcal N_x^\circ$ i
$\mathcal N_x^\circ=\mathcal N_x$.
Kładziemy $\iota(x)=\mathcal N_x$.
Ponieważ $X$ jest rozdzielająca,
różne punkty mają różne filtry otoczeń.

Pokażemy zgodność struktur na obrazie.
Jeżeli $(\iota(x),\iota(y))\in\widehat E$,
to istnieją otoczenia $A$ punktu $x$
i $B$ punktu $y$ z $A\times B\subseteq E$.
W szczególności $(x,y)\in E$.
Odwrotnie, jeśli $(x,y)\in D$
dla symetrycznego $D$ z $D^3\subseteq E$,
to $D[x]\times D[y]\subseteq E$;
obie sekcje są otoczeniami punktów,
więc $(\iota(x),\iota(y))\in\widehat E$.
To wykazuje jednostajną ciągłość
$\iota$ i jej odwrotności na obraz.

Gęstość obrazu: niech
$\mathcal F\in\widehat X$ i $E\in\mathcal U$.
Wybieramy symetryczne $D$ z $D^2\subseteq E$
oraz niepusty $A\in\mathcal F$
z $A\times A\subseteq D$.
Dla dowolnego $x\in A$ mamy
$A\times D[x]\subseteq D^2\subseteq E$,
a $D[x]\in\mathcal N_x$.
Zatem $(\mathcal F,\iota(x))\in\widehat E$.
Każde otoczenie $\mathcal F$ w $\widehat X$
spotyka $\iota(X)$.

Pozostaje zupełność. Niech $\mathcal H$
będzie filtrem Cauchy'ego na $\widehat X$.
Z gęstości $\iota(X)$ zbudujemy filtr
$\mathcal G$ na $X$ o bazie złożonej ze zbiorów
\[
 B(H,E)=\{x\in X:\exists \mathcal F\in H,\
       (\iota(x),\mathcal F)\in\widehat E\},
 \qquad H\in\mathcal H,\quad E\in\mathcal U.
\]
Są one niepuste: dla $\mathcal F\in H$
gęstość $\iota(X)$ zapewnia punkt
$\iota(x)\in\widehat E[\mathcal F]$.
Mamy
$B(H\cap H',E\cap E')\subseteq
 B(H,E)\cap B(H',E')$,
więc jest to baza właściwego filtra.

Dla $V\in\mathcal U$ wybieramy symetryczne
$D$ z $D^3\subseteq V$ i
$H\in\mathcal H$ o $H\times H\subseteq\widehat D$.
Jeżeli $x,y\in B(H,D)$, wybieramy
$\mathcal F,\mathcal G\in H$ z
$(\iota(x),\mathcal F),(\iota(y),\mathcal G)
\in\widehat D$.
Wówczas
$(\iota(x),\iota(y))\in\widehat{D^3}$,
a zgodność struktury na obrazie daje
$(x,y)\in D^3\subseteq V$.
Tak więc $\mathcal G$ jest filtrem Cauchy'ego.

Niech $\mathcal F_0=\mathcal G^\circ$
będzie jego minimalnym bliskim filtrem.
Pokażemy, że $\mathcal H$ zbiega do
$\mathcal F_0$. Dla symetrycznego
$E\in\mathcal U$ wybieramy $D$ z $D^4\subseteq E$.
Istnieją $A\in\mathcal G$ i $C\in\mathcal F_0$
z $A\times C\subseteq D$.
Ponadto możemy wybrać bazowy
$B(H,D)\subseteq A$.
Dla $\mathcal K\in H$ wybieramy
$x\in B(H,D)$ oraz
$\mathcal L\in H$ z
$(\iota(x),\mathcal L)\in\widehat D$.
Ponieważ $\mathcal H$ jest Cauchy'ego,
po zmniejszeniu $H$ na początku konstrukcji
możemy wymagać
$H\times H\subseteq\widehat D$,
więc $(\mathcal L,\mathcal K)\in\widehat D$.
Z $x\in A$ i $A\times C\subseteq D$
oraz definicji filtrów otoczeń wynika
$(\iota(x),\mathcal F_0)\in\widehat{D^2}$.
Przez złożenie otrzymujemy
$(\mathcal K,\mathcal F_0)\in\widehat{D^4}
\subseteq\widehat E$.
Zatem $H\subseteq\widehat E[\mathcal F_0]$;
ponieważ $\widehat E[\mathcal F_0]$ jest
otoczeniem, $\mathcal H\to\mathcal F_0$.
\end{proof}

\begin{uwaga}[Przestrzenie nierozdzielające]
Jeżeli $\mathcal U$ nie jest rozdzielająca,
to relacja $x\approx y$ dana warunkiem
$(x,y)\in\bigcap_{E\in\mathcal U}E$
jest relacją równoważności. Na ilorazie
otrzymuje się rozdzielającą strukturę
jednostajną, którą można uzupełnić.
Nie jest możliwe zanurzenie dowolnej
przestrzeni nierozdzielającej w rozdzielającą
przestrzeń jednostajną bez uprzedniego
utożsamienia punktów nierozróżnialnych.
\end{uwaga}

\section{Własność uniwersalna i jednoznaczność}
\label{sec:topog-uzupelnienia-uniwersalnosc}

\begin{lemat}[Rozszerzanie przez gęstość]
\label{lem:topog-uzupelnienia-rozszerzenie}
Niech $A$ będzie gęstą podprzestrzenią
jednostajną przestrzeni $Z$, a $Y$ zupełną
rozdzielającą przestrzenią jednostajną.
Każde jednostajnie ciągłe
$f:A\to Y$ ma dokładnie jedno
jednostajnie ciągłe rozszerzenie
$\overline f:Z\to Y$.
\end{lemat}
\begin{proof}
Dla $z\in Z$ oznaczmy przez $\mathcal F_z$
filtr na $A$ generowany przez
$A\cap E[z]$, gdzie $E$ przebiega
otoczenia jednostajne w $Z$.
Każdy taki zbiór jest niepusty,
gdyż $E[z]$ zawiera otwarte otoczenie
$z$, a $A$ jest gęsty. Filtr
$\mathcal F_z$ jest Cauchy'ego:
dla danego $V$ wybieramy symetryczne
$D$ z $D^2\subseteq V$;
wówczas $(D[z]\cap A)^2\subseteq V$.
Jednostajna ciągłość $f$ oznacza,
że filtr obrazów $f_*(\mathcal F_z)$
jest Cauchy'ego w $Y$. Z zupełności
zbiega do punktu $\overline f(z)$;
punkt ten jest jedyny, gdyż $Y$
jest Hausdorffa. Dla $z\in A$
filtr $\mathcal F_z$ zbiega do $z$,
więc $\overline f(z)=f(z)$.

Dowiedziemy jednostajnej ciągłości.
Ustalmy symetryczne $V$ w $Y$ i wybierzmy
symetryczne $W$ z $W^3\subseteq V$.
Z jednostajnej ciągłości $f$ istnieje
symetryczne $D$ w $Z$ takie, że
$(a,b)\in D\cap(A\times A)$ pociąga
$(f(a),f(b))\in W$.
Wybierzmy symetryczne $E$ z $E^3\subseteq D$.
Jeżeli $(z,z')\in E$, to dla
$a\in A\cap E[z]$ i $b\in A\cap E[z']$
mamy $(a,b)\in E^3\subseteq D$,
więc $(f(a),f(b))\in W$.
Przechodząc do granic filtrów w
topologii produktu $Y\times Y$,
otrzymujemy
$(\overline f(z),\overline f(z'))\in
\overline W^{\,Y\times Y}$.
Możemy na początku wybrać symetryczne $W$
takie, że $\overline W^{\,Y\times Y}\subseteq V$:
wystarczy najpierw wziąć symetryczne $T$
z $T^3\subseteq V$.
Jeśli $(y,y')\in\overline T^{\,Y\times Y}$,
to otwarte otoczenia punktów $y,y'$
zawarte odpowiednio w $T[y]$ i $T[y']$
wyznaczają otoczenie $(y,y')$ spotykające $T$.
Istnieją więc $(a,b)\in T$ z
$(y,a),(b,y')\in T$, co daje
$(y,y')\in T^3\subseteq V$.
Możemy zatem przyjąć $W=T$,
co dowodzi jednostajnej ciągłości.

Jednoznaczność: dwa ciągłe rozszerzenia
zgadzają się na gęstym $A$, więc
są równe na $Z$ dzięki Hausdorffowości $Y$.
\end{proof}

\begin{twierdzenie}[Własność uniwersalna uzupełnienia]
\label{thm:topog-uzupelnienia-uniwersalnosc}
Niech $\iota:X\to\widehat X$ będzie
uzupełnieniem rozdzielającej przestrzeni
jednostajnej $X$, skonstruowanym
w twierdzeniu~\ref{thm:topog-uzupelnienia-istnienie}.
Dla każdej zupełnej rozdzielającej
przestrzeni jednostajnej $Y$
i każdego jednostajnie ciągłego
$f:X\to Y$ istnieje dokładnie jedno
jednostajnie ciągłe $\widehat f:\widehat X\to Y$
takie, że $\widehat f\circ\iota=f$.
Ponadto dowolne dwa uzupełnienia $X$
są jednoznacznie izomorficzne jednostajnie
przez izomorfizm zachowujący obraz $X$.
\end{twierdzenie}
\begin{proof}
Zastosujmy lemat~\ref{lem:topog-uzupelnienia-rozszerzenie}
do gęstej podprzestrzeni $\iota(X)$
w $\widehat X$ i odwzorowania
$f\circ\iota^{-1}$. Otrzymujemy
żądane rozszerzenie.

Dla dwóch uzupełnień $X_1,X_2$
własność uniwersalna daje jednostajnie
ciągłe odwzorowania
$T:X_1\to X_2$ i $S:X_2\to X_1$
zgodne z odpowiednimi zanurzeniami.
Złożenia $S\circ T$ i $T\circ S$
są identycznościami na gęstych
obrazach $X$; z jednoznaczności
rozszerzenia są identycznościami
na $X_1$ i $X_2$. Tak więc
$T$ jest izomorfizmem jednostajnym.
\end{proof}

\section{Domknięte podprzestrzenie i przykłady}
\label{sec:topog-uzupelnienia-przyklady}

\begin{twierdzenie}[Zupełność podprzestrzeni]
\label{thm:topog-uzupelnienia-podprzestrzenie}
Niech $X$ będzie zupełną rozdzielającą
przestrzenią jednostajną i $A\subseteq X$
z odziedziczoną strukturą jednostajną.
Wówczas $A$ jest zupełna wtedy i tylko wtedy,
gdy jest domknięta w $X$.
\end{twierdzenie}
\begin{proof}
Jeżeli $A$ jest domknięta, to każdy filtr
Cauchy'ego na $A$, traktowany przez
obrazy zbiorów jako filtr na $X$,
jest Cauchy'ego. Ma granicę $x\in X$.
Ponieważ filtr ma bazę złożoną ze
zbiorów zawartych w $A$, granica należy
do $\overline A=A$.
Stąd filtr zbiega w $A$.

Odwrotnie, niech $A$ będzie zupełna
i $x\in\overline A$. Filtr na $A$
generowany przez $A\cap E[x]$ jest
Cauchy'ego, jak w dowodzie
lematu~\ref{lem:topog-uzupelnienia-rozszerzenie}.
Zbiega więc do pewnego $a\in A$.
Jako filtr na $X$ zbiega zarazem do $x$.
Hausdorffowość $X$ daje $a=x$;
stąd $A$ jest domknięta.
\end{proof}

\begin{przyklad}[Dwa uzupełnienia tej samej topologii]
\label{ex:topog-uzupelnienia-dwie-metryki}
Uzupełnienie $(\R,|\cdot|)$ jest równe
$\R$ z tą samą metryką. Z kolei dla
$e(x,y)=|\arctan x-\arctan y|$
izometria $\arctan$ identyfikuje $\R$
z otwartym przedziałem
$(-\pi/2,\pi/2)$, którego uzupełnieniem
jest $[-\pi/2,\pi/2]$.
Uzupełnienia te nie są nawet
homeomorficzne: pierwsze jest niezwartą
przestrzenią, a drugie zwartą.
\end{przyklad}

\begin{przyklad}[Produkt i zupełność]
\label{ex:topog-uzupelnienia-produkt}
Jeżeli każda $(X_i,\mathcal U_i)$ jest
zupełna, to produkt $\prod_{i\in I}X_i$
ze strukturą produktu z rozdziału~12
jest zupełny. Istotnie, filtr Cauchy'ego
$\mathcal F$ w produkcie przenosi się
przez każde rzutowanie na filtr Cauchy'ego
w $X_i$. Wybierając jego granicę $x_i$
dla każdego $i$ (dla pustego produktu
otrzymujemy jedyny punkt), dostajemy
$x=(x_i)$, do którego zbiega
$\mathcal F$: każde otwarte otoczenie
$x$ zawiera cylinder oparty na skończenie
wielu otoczeniach $x_i$, a przeciwobrazy
tych otoczeń należą do $\mathcal F$.
Gdy którykolwiek czynnik jest pusty,
produkt jest pusty i zupełny w sensie
definicji~\ref{def:topog-uzupelnienia-zupelnosc}.
\end{przyklad}

\begin{uwaga}[Kierunek dalszej pracy]
Każda zwarta przestrzeń jednostajna jest
zupełna. Rzeczywiście, filtr Cauchy'ego
$\mathcal F$ w zwartej przestrzeni $X$
ma punkt skupienia $x$: rodzina
$\{\overline A:A\in\mathcal F\}$ ma własność
skończonych przecięć, więc jej przecięcie
jest niepuste. Ustalmy $E\in\mathcal U$
i symetryczne $D$ z $D^2\subseteq E$.
Istnieje $A\in\mathcal F$ z
$A\times A\subseteq D$. Ponieważ
$x\in\overline A$ i $D[x]$ jest otoczeniem
$x$, istnieje $a\in A\cap D[x]$.
Dla każdego $b\in A$ mamy
$(x,a),(a,b)\in D$, a zatem
$b\in E[x]$. Stąd $E[x]\in\mathcal F$,
czyli $\mathcal F\to x$.
Konstrukcja uzupełnienia
nie jest jednak uzwarceniem:
uzupełnienie nie musi być zwarte,
a jego funkcją jest domknięcie braków
w filtrach Cauchy'ego, nie dodanie granic
wszystkim filtrom. Powiązania z
uzwarceniami rozważymy oddzielnie.
\end{uwaga}
